\section{Recomendaciones}
    \begin{enumerate}
        \item El sector explica el 84\% de la variabilidad del consumo,
        por lo que mezclar los tres sectores en una sola cifra puede ocultar
        diferencias relevantes para pronósticos, tarifas y metas de ahorro.
        
        \vspace{5pt}
        
        \item En Residencial existe una relación fuerte ($r=0{,}595$) entre
        temperatura y consumo, mientras que en el modelo agregado parece nula
        ($r=0{,}063$), por ello, la temperatura constituye una variable útil
        para anticipar picos de demanda.
        
        \vspace{5pt}
        
        \item La dispersión de Industrial (195 kWh) es casi cinco veces la de
        Residencial (41 kWh), diferencia que debe considerarse en los contratos
        de suministro y márgenes de seguridad.
        
        \vspace{5pt}
        
        \item Un valor $p$ pequeño no mide la magnitud del efecto, ni uno
        grande demuestra su ausencia, por ende, es necesario complementar la
        significancia con $d$, $\omega^2$ y la potencia estadística.
        
        \vspace{5pt}
        
        \item Una conclusión de ``no hay relación'' debe acompañarse de una
        visualización desagregada que, para este caso, la Figura \ref{fig_plotly}
        permitió identificar la relación que el coeficiente global ocultaba.
    \end{enumerate}

    \vspace{10pt}

    El código completo, las tablas y las figuras interactivas en HTML
    están disponibles en el repositorio del proyecto:

    \begin{center}
        \begin{tcolorbox}[
            arc=0pt, boxrule=0.8pt,
            colback=white, colframe=green!45!black,
            colbacktitle=green!45!black, coltitle=white,
            fonttitle=\bfseries,
            title=Repositorio del proyecto en GitHub,
            halign=center, halign title=center,
            width=0.92\linewidth
        ]
            \small
            \url{https://github.com/RubianoAndy/hypothesis_visualizations}
        \end{tcolorbox}
    \end{center}
